\subsection{Descripción del dataset}
Se utiliza el archivo \path{WA_Fn-UseC_-Telco-Customer-Churn.csv} de la publicación Telco 
Customer Churn de BlastChar en Kaggle~\cite{dataset}. El archivo contiene información de 
7.043 clientes y 21 columnas, incluyendo variables demográficas, de servicio, de cuenta 
y de facturación. Cada fila corresponde a un cliente. Las columnas incluyen un identificador, 
la variable objetivo \var{Churn} y 19 características que podrían utilizarse para el modelado. 
De estas, tres son numéricas y 16 categóricas. 

\var{SeniorCitizen} es categórica binaria aunque se almacene como entero. \var{TotalCharges} inicialmente 
se lee como texto y se convierte a número. La Tabla~\ref{tab:variables} resume las columnas y sus tipos de variables.

\begin{table*}[t]
    \caption{Variables del dataset y sus tipos}
    \label{tab:variables}
    \centering\footnotesize

    \begin{tabularx}{\textwidth}{
        @{}p{1.8cm}
        >{\raggedright\arraybackslash}X
        >{\raggedright\arraybackslash}p{3.2cm}@{}
    }
        \toprule
        Grupo & Variables & Tipo o descripción \\
        \midrule
        Identificador & \var{customerID}
        & Texto, único por cliente \\

        Objetivo & \var{Churn}
        & Categoría binaria (No/Yes) \\

        Perfil & \var{gender}, \var{SeniorCitizen},
        \var{Partner}, \var{Dependents}
        & Categorías binarias \\

        Antigüedad & \var{tenure}
        & Número entero de meses \\

        Telefonía & \var{PhoneService}, \var{MultipleLines}
        & Categóricas \\

        Internet & \var{InternetService}
        & Categórica (DSL, Fiber optic o No) \\

        Adicionales & \var{OnlineSecurity}, \var{OnlineBackup},
        \var{DeviceProtection}, \var{TechSupport},
        \var{StreamingTV}, \var{StreamingMovies}
        & Categóricas \\

        Cuenta & \var{Contract}, \var{PaperlessBilling},
        \var{PaymentMethod}
        & Categóricas \\

        Cargos & \var{MonthlyCharges}, \var{TotalCharges}
        & Numéricas, cargo mensual y acumulado \\
        \bottomrule
    \end{tabularx}
\end{table*}